\documentclass[11pt]{article}

\usepackage[a4paper,margin=1in]{geometry}
\usepackage{amsmath}
\usepackage{amssymb}
\usepackage{booktabs}
\usepackage{enumitem}
\usepackage[hidelinks]{hyperref}

\emergencystretch=2em
\Urlmuskip=0mu plus 2mu

\title{Planar Tensor Networks, Braiding, and Interface Invariants}
\author{}
\date{\today}

\newcommand{\catC}{\mathcal{C}}
\newcommand{\Hom}{\operatorname{Hom}}
\newcommand{\id}{\operatorname{id}}
\newcommand{\Space}{\mathsf{Space}}
\newcommand{\Dual}{\mathsf{Dual}}
\newcommand{\Leg}{\mathsf{Leg}}
\newcommand{\Tensor}{\mathsf{Tensor}}
\newcommand{\Network}{\mathsf{Network}}
\newcommand{\Plan}{\mathsf{Plan}}
\newcommand{\cod}{\operatorname{cod}}
\newcommand{\dom}{\operatorname{dom}}

\begin{document}

\maketitle

\begin{abstract}
This note records the design consequences of supporting non-Abelian
symmetries, fermionic signs, and general braiding in a tensor-network
library.  The main point is simple: once braiding or nontrivial
recoupling is present, a tensor expression is not determined by the
set of labels on its legs.  It is determined by an ordered and
possibly planar diagram.  A robust interface therefore needs explicit
objects for spaces, dual spaces, ordered legs, fusion trees, planar
network geometry, and contraction plans.  Mutable labels are semantic
compatibility tags on space values, but they cannot replace canonical
axis positions or the structural data of tensor legs.
\end{abstract}

\tableofcontents

\section{Purpose}

The immediate motivation is the failure mode of Cytnx-style tensor
labels.  In Cytnx, labels are persistent mutable strings attached to
tensor legs, and many contractions use those labels to decide which
legs to contract.  This can be convenient for dense tensors, but it is
not a sound foundation for tensor categories with nontrivial
associators, fermionic signs, or braiding.

The target design here is closer to what a future \texttt{uni20}
interface should enforce.  It is not a proposal to reproduce the
Cytnx C++ object model.  In particular:

\begin{itemize}[nosep]
  \item dense arrays should be ordinary array objects, not the central
        tensor-network abstraction;
  \item tensor-network correctness should be expressed in terms of
        spaces, dual spaces, ordered legs, morphisms, and diagrams;
  \item high-level notation may look like einsum, but it must compile
        to a checked contraction plan;
  \item labels are useful for compatibility and diagnostics, but they
        must not select or reorder axes.
\end{itemize}

\section{Minimal Categorical Model}

A tensor-network library with symmetries is not merely an ndarray
library with labels.  It is manipulating morphisms in a monoidal
category \(\catC\).  The objects of \(\catC\) are spaces, and tensors
are maps between tensor products of spaces.

For a tensor \(A\), write
\[
  A : X_1 \otimes X_2 \otimes \cdots \otimes X_m
      \longrightarrow
      Y_1 \otimes Y_2 \otimes \cdots \otimes Y_n .
\]
Equivalently,
\[
  A \in \Hom_{\catC}
       (X_1 \otimes \cdots \otimes X_m,
        Y_1 \otimes \cdots \otimes Y_n).
\]

The distinction between domain and codomain is not optional.  A
contraction pairs a space with a compatible dual space.  For ordinary
finite-dimensional vector spaces, this distinction can be flattened
away by choosing bases, but tensor-network code with symmetries should
not erase it.

\subsection{Strict and Symmetric Cases}

Dense tensors over ordinary vector spaces live in a symmetric monoidal
category.  In this setting:

\begin{itemize}[nosep]
  \item associators can usually be treated as identities;
  \item swapping two factors is represented by an ordinary permutation;
  \item there is no sign, \(R\)-matrix, or recoupling coefficient;
  \item many illegal-looking relabellings accidentally produce the
        intended dense numerical answer.
\end{itemize}

This is the misleading special case.  It trains users to believe that
``if the labels are right, the order does not matter.''  That is not a
valid rule in the general setting.

\subsection{Non-Abelian and Braided Cases}

For non-Abelian symmetries, a basis for a multi-leg tensor depends on
the fusion tree.  For example,
\[
  ((a \otimes b) \otimes c) \to d
\]
and
\[
  (a \otimes (b \otimes c)) \to d
\]
are related by an \(F\)-move.  They are not the same basis with
different leg names.

For braided tensor categories, exchanging two objects applies an
\(R\)-move:
\[
  R_{a,b}: a \otimes b \to b \otimes a .
\]
For fermionic tensors this may reduce to parity signs, but the same
interface lesson applies: a swap is a mathematical operation, not a
metadata edit.

\section{What ``Planar'' Means}

In this note, a planar tensor network is not just a graph.  It is a
graph with enough embedding data to say which tensor legs are adjacent
and whether a proposed contraction crosses another strand.

At minimum, a planar network needs:

\begin{itemize}
  \item vertices, each representing a tensor or elementary morphism;
  \item half-edges attached to ordered tensor legs;
  \item internal edges pairing compatible half-edges;
  \item external boundary legs, in a specified boundary order;
  \item a cyclic order of half-edges around every vertex;
  \item enough embedding information to decide whether adding or
        reordering an edge introduces a crossing;
  \item explicit crossing or braiding nodes when crossings are intended.
\end{itemize}

For a dense bosonic tensor network, most of this can be ignored.  For a
braided category, it is part of the expression.

\subsection{Local Order vs Global Planarity}

A single tensor can store a local ordered list of legs:
\[
  [\ell_0,\ell_1,\ldots,\ell_{r-1}].
\]
That is necessary, but not sufficient.  It says the order around that
one tensor.  It does not fully determine the embedding of a multi-tensor
network.

The global network must additionally know how tensor vertices are
placed relative to each other, or at least an equivalent combinatorial
embedding.  Without this information, the library cannot distinguish
between these two cases:

\[
  \text{two non-crossing contractions}
  \qquad\text{vs.}\qquad
  \text{the same pairings with one crossing}.
\]

In a symmetric category those two diagrams are equivalent after an
ordinary permutation.  In a braided category they differ by \(R\)-moves.

\subsection{Why a Tensor Alone Cannot Know Whether a Braid Is Needed}

Whether a braid is needed is not a property of one tensor in isolation.
It depends on the larger expression.

For example, suppose tensor \(A\) has outgoing boundary order
\[
  [a,b]
\]
and tensor \(B\) has incoming boundary order
\[
  [a^\ast,b^\ast].
\]
Contracting \(a\) with \(a^\ast\) and \(b\) with \(b^\ast\) may be
planar in one embedding.  Contracting \(a\) with \(b^\ast\) and \(b\)
with \(a^\ast\) requires a crossing unless the tensors are embedded
differently or explicit braids have been inserted.

The individual tensors only know their local boundary orders.  The
network or contraction planner must decide whether the pairings are
planar.

\section{Three-Layer Design}

A robust implementation should separate three layers.

\subsection{Layer 1: Tensor Objects}

A tensor object stores local categorical data:

\begin{itemize}
  \item an ordered domain list;
  \item an ordered codomain list;
  \item the corresponding spaces and dual spaces;
  \item semantic string labels attached to the concrete space values;
  \item symmetry-sector data and degeneracy data;
  \item fusion-tree data, if the category is non-Abelian;
  \item reduced numerical blocks.
\end{itemize}

This layer can answer questions such as:

\begin{itemize}[nosep]
  \item What is the \(i\)-th leg?
  \item Is this leg a concrete space or \(\Dual\langle S\rangle\)?
  \item What sectors live on this leg?
  \item What fusion tree defines the basis?
  \item Does this tensor have the same boundary type as another tensor?
\end{itemize}

It cannot, by itself, decide all global planarity questions.

\subsection{Layer 2: Planar Network Objects}

A network object stores a diagram:

\begin{itemize}
  \item tensor vertices;
  \item internal pairings;
  \item external boundary order;
  \item cyclic orders around vertices;
  \item explicit braid/crossing vertices if present;
  \item optional geometric coordinates for display, not as the
        authority on topology.
\end{itemize}

The network object is the right place to ask:

\begin{itemize}[nosep]
  \item Is this proposed contraction planar?
  \item Would contracting these two legs cross an existing edge?
  \item What is the boundary order after contraction?
  \item Is an explicit braid required?
  \item Can this diagram be simplified by isotopy?
\end{itemize}

\subsection{Layer 3: Contraction Plans}

A contraction plan is an executable sequence of local operations.  It
may be produced from a network after checking the diagram.

The plan records:

\begin{itemize}
  \item which local contractions are performed;
  \item which tensor products are formed;
  \item which transposes are merely cyclic reorders;
  \item which \(F\)-moves are inserted;
  \item which \(R\)-moves are inserted;
  \item which traces include twists or supertraces;
  \item the expected final boundary order and spaces.
\end{itemize}

This is the layer that can optimize.  The tensor object should not
silently optimize away categorical distinctions.

\section{Planar Invariants}

The following invariants are the core of the design.  They are the
rules that should remain true after every operation.

\subsection{Boundary Invariant}

Every tensor and every network has a well-defined ordered boundary.
For a tensor this is the ordered domain and codomain.  For a network it
is the ordered list of uncontracted external legs.

An operation must specify the boundary it produces.  It is not enough
to produce numerically shaped data.

\subsection{Dual Compatibility Invariant}

A contraction can only pair compatible dual spaces.  In the simplest
case this means \(V\) with \(V^\ast\).  With symmetries it also means
compatible charge sectors, degeneracies, and orientations.

Equal total dimension is not sufficient.  Equal label string is not
sufficient.

\subsection{Space Compatibility Invariant}

Two spaces may have the same dimension and the same charge sectors but
still represent different positions in a tensor network.  For example,
two MPS bond spaces at different cuts may both have dimension \(32\).
They should not be interchangeable unless their complete space values
are equal or the operation supplies an explicit structural
isomorphism.

This is where labels are useful when they belong to space values rather
than forming a second axis-addressing system.  A finite MPS can label
its bond spaces
\[
  b_0,b_1,\ldots,b_L,
\]
and then exact compatibility rejects a wrong bond even in the dense
case.  Exact equality includes immutable structure and the mutable
label.  There is no separate globally generated space identifier.

\subsection{Cyclic Order Invariant}

At every tensor vertex, the cyclic or linear order of legs is
structural data.  Reordering legs changes the tensor expression unless
the operation is explicitly known to be valid.

Dense tensors may implement this as a cheap axis permutation.  Braided
or non-Abelian tensors cannot treat it as mere bookkeeping.

\subsection{No Implicit Crossing Invariant}

No operation may introduce a crossing silently.

If a proposed contraction is non-planar, the operation should fail with
a diagnostic such as:

\begin{quote}
This contraction is not planar in the current diagram.  Insert an
explicit braid, or use a contraction planner that is allowed to insert
braids.
\end{quote}

The important point is that the user learns what mathematical operation
is missing.

\subsection{Fusion-Tree Invariant}

For non-Abelian categories, every tensor basis is defined relative to a
fusion tree or an equivalent canonical representation.

If an operation changes the parenthesization of a tensor product, it
must either:

\begin{itemize}[nosep]
  \item prove that no recoupling is needed;
  \item apply the appropriate \(F\)-move;
  \item or fail.
\end{itemize}

Changing labels must never be able to bypass this.

\subsection{Braid Accounting Invariant}

If two strands are exchanged in a braided category, the corresponding
\(R\)-move must appear in the tensor expression, either as an explicit
operation or as an explicit node in the contraction plan.

The library should distinguish:

\begin{itemize}[nosep]
  \item a permutation in a symmetric category;
  \item a fermionic swap with parity signs;
  \item a genuine anyonic braid with an \(R\)-matrix;
  \item a display-only reorder.
\end{itemize}

\subsection{Trace and Twist Invariant}

Closing a line into a trace is not always the same operation as a
matrix trace.  In fermionic or braided settings, trace-like operations
may require twists, signs, or categorical traces.

The interface should therefore distinguish ordinary dense traces from
categorical traces.  A planner should not silently replace one with
the other.

\subsection{Decomposition Invariant}

SVD, QR, LQ, polar decomposition, and related operations are not just
matrix factorizations.  In a tensor-network context they create new
intermediate spaces.

For an SVD of
\[
  A : X \to Y,
\]
the decomposition
\[
  A = U S V
\]
creates an intermediate space \(B\).  The singular-value object \(S\)
is an endomorphism of \(B\):
\[
  S : B \to B.
\]
Thus \(S\) naturally has equal copies of the same labeled space value
on its two sides.  A
global rule that tensor labels must be unique is incompatible with this
mathematical structure.

\subsection{Labels Are Not Axis Operations}

String labels express intended leg compatibility and participate in
exact space equality.  Names, colors, and diagram coordinates may also
help humans.  None of them chooses an axis or changes sector structure.

Changing a label may deliberately make two leg values equal or unequal.
Changing leg order, structural decomposition, or braiding requires a
separate checked operation.

\section{Operation Taxonomy}

The interface should split operations that dense tensor libraries often
conflate.

\begin{center}
\begin{tabular}{lll}
\toprule
Operation & Dense interpretation & General categorical interpretation \\
\midrule
\texttt{relabel} & rename a leg space & compatibility metadata only \\
\texttt{permute} & reorder axes & may require braid or recoupling \\
\texttt{transpose} & swap matrix axes & cyclic/domain-codomain move, if valid \\
\texttt{reshape} & reinterpret shape & usually not meaningful for spaces \\
\texttt{contract} & sum over axes & compose compatible morphisms \\
\texttt{trace} & sum diagonal & categorical trace may need twist/sign \\
\texttt{svd} & matrix SVD & factorization creating a new bond space \\
\texttt{get\_block} & array slice & reduced block in a chosen fusion basis \\
\bottomrule
\end{tabular}
\end{center}

\subsection{Reorder Operations}

At least four operations should be distinct:

\begin{description}
  \item[\texttt{relabel}] Change semantic leg labels only.
  \item[\texttt{repartition}] Move the domain/codomain boundary without
        changing the cyclic order, when categorically valid.
  \item[\texttt{transpose}] Perform a cyclic reorder that introduces no
        crossing.
  \item[\texttt{braid}] Exchange adjacent strands using an \(R\)-move.
\end{description}

Dense libraries collapse these into one or two functions.  That is
exactly what should be avoided in a category-aware tensor-network
interface.

\section{TensorKit.jl as a Useful Reference Point}

TensorKit.jl provides a useful example of this separation.  I am not
claiming uni20 should copy its interface exactly, but its design
validates the main point of this note: once braiding matters, ordinary
implicit tensor notation is not enough.

The TensorKit documentation describes tensors as \texttt{TensorMap}
objects, i.e. maps from a domain to a codomain.  It also has explicit
fusion trees, with uncoupled sectors, coupled sectors, duality flags,
inner fusion channels, and vertex labels.  This is precisely the kind
of information that a Cytnx-style leg label cannot encode.

\subsection{Fusion Trees}

TensorKit's fusion-tree data stores the basis structure of a symmetric
tensor.  The important lesson is that block sparsity is not just a
dictionary keyed by quantum numbers.  It is block data together with a
choice of fusion basis.

For uni20, this suggests:

\begin{itemize}[nosep]
  \item tensor blocks should be reduced data;
  \item fusion paths should be explicit or canonical;
  \item recoupling should be an operation on the basis, not an
        accidental consequence of permuting axes.
\end{itemize}

\subsection{TensorKit's Hierarchy of Reordering Operations}

TensorKit distinguishes several index manipulations:

\begin{description}
  \item[\texttt{braid}] The most general operation.  It can represent
        crossings and needs information about over/under order.
  \item[\texttt{permute}] A restricted operation for sector types with
        symmetric braiding, where over/under crossings are equivalent.
  \item[\texttt{transpose}] Cyclic permutations that do not introduce
        crossings.
  \item[\texttt{repartition}] Moving the boundary between domain and
        codomain without arbitrary reordering.
\end{description}

This is a useful vocabulary.  It explicitly prevents the common dense
tensor mistake of treating every leg reorder as an axis permutation.

\subsection{\texttt{@tensor}, \texttt{@planar}, and \texttt{@plansor}}

TensorKit's contraction macros make the planar issue concrete.

For bosonic or symmetric cases, ordinary tensor notation can be enough.
TensorKit's \texttt{@tensor} macro handles implicit contractions in
that setting.

For anyonic or braided cases, the same implicit notation is not
sufficient.  TensorKit provides a \texttt{@planar} macro where the
expression must describe a planar diagram.  Crossings are not silently
inferred; they are represented explicitly by a reserved braiding tensor
such as \(\tau\), with \(\tau'\) for the inverse crossing.

TensorKit also provides \texttt{@plansor}, which dispatches to the
ordinary tensor machinery in bosonic cases and to planar machinery when
the category requires it.

The design lesson is strong: there are two different languages.

\begin{itemize}
  \item A non-planar indexed expression language, acceptable only when
        the category is symmetric enough.
  \item A planar diagram language, required when crossings have
        mathematical content.
\end{itemize}

\section{A Possible uni20 Interface Shape}

The following sketches are illustrative, not final API.

\subsection{Spaces and Legs}

\begin{verbatim}
phys = LocalSpace(SU2, [j(1,2)], "SpinHalf")
b0   = BlockSpace(SU2, [(j(0), 1)], "bond-0")
b1   = BlockSpace(SU2, [(j(0), 1), (j(1), 1)], "bond-1")

A = Tensor(
    codomain=[b1],
    domain=[b0, phys],
    fusion_tree=...
)
\end{verbatim}

The key point is that the tensor is a morphism.  Its boundary is
ordered.  Labels such as \texttt{"bond-0"} are attached to spaces for
diagnostics and compatibility checks.

\subsection{Contraction by Explicit Legs}

\begin{verbatim}
C = contract(A, B, pairs=[(A.codomain[0], B.domain[0])])
\end{verbatim}

This operation checks:

\begin{itemize}[nosep]
  \item dual compatibility;
  \item exact space equality or an allowed structural isomorphism;
  \item sector compatibility;
  \item fusion-tree compatibility;
  \item planarity in the current network context.
\end{itemize}

\subsection{Network-Level Planar Contraction}

\begin{verbatim}
N = PlanarNetwork()
N.add_tensor("A", A)
N.add_tensor("B", B)
N.connect(("A", 0), ("B", 1))

plan = N.compile(boundary=[...])
C = plan.evaluate()
\end{verbatim}

If a crossing is needed, compilation fails unless it is present in the
diagram or explicitly authorized:

\begin{verbatim}
N.braid(("A", 1), ("B", 0), over=True)
\end{verbatim}

\subsection{Einsum-Like Syntax as a Front End}

An einsum-like expression can still exist:

\begin{verbatim}
C = einsum("lir,rjs->lijs", A, B)
\end{verbatim}

but this should be a notation parser.  It should produce a contraction
plan and then check that the local labels agree with the tensors'
actual spaces.  The labels in the expression should not mutate the
tensors.

\section{Diagnostics}

The error messages matter.  The library should tell users which
categorical condition failed.

Examples:

\begin{quote}
Cannot contract leg 2 of \(A\) with leg 0 of \(B\): spaces have equal
dimension and sectors but different labels, \texttt{bond-4} vs
\texttt{bond-5}.
\end{quote}

\begin{quote}
Cannot contract these legs in the current planar diagram: the pairing
would cross an existing strand.  Insert an explicit braid or use a
planner with braiding enabled.
\end{quote}

\begin{quote}
Cannot reorder legs by metadata relabel.  This reorder changes the
fusion tree.  Apply an explicit \(F\)-move or call \texttt{recouple}.
\end{quote}

\begin{quote}
Cannot use ordinary trace on this fermionic/braided tensor.  Use the
categorical trace operation and specify the required twist convention.
\end{quote}

These diagnostics are not just niceties.  They are part of the
correctness model.  A useful tensor-network library should make the
wrong operation fail at the point where the wrong assumption is made.

\section{Testing the Invariants}

The invariants above should be testable.

\subsection{Dense Space Compatibility Tests}

Even with no symmetry, create two \texttt{DenseSpace} values with the
same extent but different labels:
\[
  V = \mathbb{C}^{32}, \qquad W = \mathbb{C}^{32}.
\]
Contraction of \(V\) with \(W^\ast\) should fail unless an explicit
isomorphism is supplied.  This catches many MPS bond-index mistakes
that dense tensor libraries silently allow.

\subsection{SVD Bond Creation Tests}

SVD should create a new bond space \(B\).  The singular-value object
should be an endomorphism \(B \to B\).  Tests should verify that the
two sides of \(S\) carry equal bond-space values and that \(U S V\)
reconstructs the original tensor with the expected boundary.

\subsection{Relabel Compatibility Tests}

Changing a label changes exact compatibility but not structural
compatibility.  Tests should verify both relations and should verify
that a network-level bond-label operation updates both endpoints.

\subsection{Braiding Tests}

In a category with nontrivial \(R\)-symbols, compare:

\begin{itemize}[nosep]
  \item a planar contraction with no crossing;
  \item the same pairings with one explicit braid;
  \item a rejected attempt to produce the second result by relabelling
        or ordinary permutation.
\end{itemize}

The first two may both be legal, but they should generally produce
different morphisms.  The third should fail.

\subsection{Fusion Recoupling Tests}

Construct tensors in two different fusion bases,
\[
  ((a b)c)d
  \qquad\text{and}\qquad
  (a(b c))d.
\]
Verify that direct comparison or contraction fails unless an explicit
\(F\)-move or canonical recoupling is applied.

\section{Consequences for Cytnx Compatibility}

A Cytnx-compatible migration layer can preserve some surface spelling,
but it should not preserve unsafe semantics.

\begin{itemize}
  \item Persistent labels can be read as semantic leg labels.
  \item Label-based contractions can be parsed into explicit leg
        contractions, then checked.
  \item Duplicate labels should not be globally forbidden, because
        endomorphisms naturally have the same space on both sides.
  \item Relabelling can change exact compatibility, but it must not
        perform a permutation, braid, recoupling, or structural
        conversion.
  \item Rowrank-style positional decomposition should be replaced by
        explicit left/right leg groups.
  \item Block access should be expert-level and described as reduced
        basis access, not ordinary ndarray slicing.
\end{itemize}

This means a fully bug-compatible shim is neither possible nor
desirable.  The useful compatibility target is:

\begin{quote}
Old scripts either run with the intended semantics, or fail with a
diagnostic that tells the user what mathematical structure was missing.
\end{quote}

\section{Conclusion}

Planarity is not decoration.  In a braided or non-Abelian tensor
category, the diagram is part of the expression.  A tensor object knows
its local ordered boundary, but the network or planner must know the
global planar embedding.  Any interface that lets users change axis
order or erase a braid by mutating labels has already lost the
essential correctness property.

The practical design rule is:

\begin{quote}
Use labels for intended leg compatibility and diagnostics.  Use
positions, space structure, duals, ordered legs, fusion trees, and
planar diagrams for the remaining semantics.
\end{quote}

\begin{thebibliography}{9}

\bibitem{TensorKitDocs}
TensorKit.jl documentation.
\newblock \url{https://quantumkithub.github.io/TensorKit.jl/stable/}

\bibitem{TensorKitFusionTrees}
TensorKit.jl manual: fusion trees.
\newblock \url{https://quantumkithub.github.io/TensorKit.jl/stable/man/fusiontrees/}

\bibitem{TensorKitIndexManipulations}
TensorKit.jl manual: index manipulations.
\newblock \url{https://quantumkithub.github.io/TensorKit.jl/stable/man/indexmanipulations/}

\bibitem{TensorKitContractions}
TensorKit.jl manual: contractions, including \texttt{@tensor},
\texttt{@planar}, and \texttt{@plansor}.
\newblock \url{https://quantumkithub.github.io/TensorKit.jl/stable/man/contractions/}

\end{thebibliography}

\end{document}
